\section*{Chapitre \ref{ch:b1:spherical-astronomy} --- Astronomie sphérique}

\begin{solution}{exo:b1:spherical-astronomy:1}
La \omterm{def:b1:spherical-astronomy:horizontal}{hauteur} du pôle est égale à la \omterm{def:b1:spherical-astronomy:geographic}{latitude}~: \ang{52}. Une étoile ne se couche
jamais si $\delta > \ang{90} - \ang{52} = \ang{38}$, et ne se lève jamais si
$\delta < \ang{-38}$.
\end{solution}

\begin{solution}{exo:b1:spherical-astronomy:2}
$\qty{6.75}{h}\times\ang{15} = \ang{101.25}$. $H = \mathrm{LST} - \alpha =
\qty{10}{h} - \qty{6}{h}\,\qty{45}{min} = \qty{3}{h}\,\qty{15}{min}$, soit
\ang{48.75}. L'\omterm{def:b1:spherical-astronomy:equatorial}{angle horaire} est positif~: l'étoile est à l'ouest du méridien,
qu'elle a traversé il y a trois heures sidérales et quart.
\end{solution}

\begin{solution}{exo:b1:spherical-astronomy:3}
$h = \ang{90} - \varphi + \delta_\odot$~: \ang{73.44} au solstice de juin,
\ang{50} aux équinoxes, \ang{26.56} au solstice de décembre.
\end{solution}

\begin{solution}{exo:b1:spherical-astronomy:4}
1900~: divisible par 100 mais pas par 400, ce n'est pas une année bissextile.
2000~: divisible par 400, année bissextile. 2100~: pas bissextile. L'année
julienne dépasse l'\omterm{def:b1:spherical-astronomy:tropical-year}{année tropique} de $365.25 - 365.2422 = 0.0078$ jour, soit
\num{0.78} jour par siècle~; l'année grégorienne, de $0.0003$ jour, soit
\num{0.03} jour par siècle.
\end{solution}

\begin{solution}{exo:b1:spherical-astronomy:5}
$\cos\Delta = \sin^2\ang{45} + \cos^2\ang{45}\cos\ang{90} = 0.5$, donc
$\Delta = \ang{60}$ et la distance vaut $6371\times\pi/3 = \qty{6672}{km}$.
Le long du parallèle, \omterm{def:b1:spherical-astronomy:spherical-triangle}{petit cercle} de rayon $R\cos\ang{45}$, la distance vaut
$6371\times\cos\ang{45}\times\pi/2 = \qty{7076}{km}$. Le \omterm{def:b1:spherical-astronomy:spherical-triangle}{grand cercle} est plus
court d'environ \qty{404}{km}.
\end{solution}

\begin{solution}{exo:b1:spherical-astronomy:6}
$\sin h = \sin\ang{48}\sin\ang{20} + \cos\ang{48}\cos\ang{20}\cos\ang{30}
= 0.2542 + 0.5446 = 0.7987$, $h = \ang{53.0}$.
$\cos\mathrm{A} = (\sin\ang{20} - \sin\ang{48}\sin\ang{53.0})/(\cos\ang{48}
\cos\ang{53.0}) = (0.3420 - 0.5935)/0.4027 = -0.6246$, $\arccos = \ang{128.7}$~;
l'étoile est à l'ouest du méridien ($H > 0$), donc $\mathrm{A} = \ang{360} -
\ang{128.7} = \ang{231.3}$~: vers le sud-ouest, à \ang{53} au-dessus de l'horizon.
\end{solution}

\begin{solution}{exo:b1:spherical-astronomy:7}
$\cos H_0 = -\tan\ang{60}\tan\ang{23.44} = -1.7321\times0.4335 = -0.7510$,
$H_0 = \ang{138.67}$, durée du jour $2\times138.67/15 = \qty{18.49}{h}$, soit
environ \qty{18}{h}~\qty{29}{min}, comme le montre la courbe de \ang{60}~N au
solstice de juin.
\end{solution}

\begin{solution}{exo:b1:spherical-astronomy:8}
Le passage avance de \qty{3.93}{min} par jour, donc de $30\times3.93 =
\qty{118}{min}$ en 30 jours~: l'étoile traverse le méridien vers 20~h~02. Pour
qu'elle le traverse à midi, il faut qu'il avance de \qty{10}{h} = \qty{600}{min},
ce qui prend $600/3.93 \approx 153$ jours.
\end{solution}

\begin{solution}{exo:b1:spherical-astronomy:9}
(a) À l'équinoxe, la \omterm{def:b1:spherical-astronomy:horizontal}{hauteur} à midi est de \ang{45}~: $I = 1361\times\sin\ang{45}
= \qty{962}{W.m^{-2}}$. (b) Au solstice de juin, $h = \ang{68.44}$ et
$(a/r)^2 = (149.598/152.100)^2 = 0.9674$~:
$I = 1361\times0.9674\times\sin\ang{68.44} = \qty{1224}{W.m^{-2}}$. Le Soleil
plus haut compense largement l'éloignement.
\end{solution}

\begin{solution}{exo:b1:spherical-astronomy:10}
$\sin\delta_\odot = \sin\ang{23.44}\sin\ang{60} = 0.3445$,
$\delta_\odot = \ang{20.15}$. $\tan\alpha_\odot = \cos\ang{23.44}\tan\ang{60}
= 1.5890$, $\alpha_\odot = \ang{57.82}$ (dans le même quadrant que
$\lambda_\odot$). L'\omterm{def:b1:spherical-astronomy:equatorial}{ascension droite} est en retard de \ang{2.18}, soit
$2.18\times4 \approx
\qty{8.7}{min}$ de temps. Un Soleil qui se déplacerait uniformément le long de
l'\omterm{def:b1:spherical-astronomy:ecliptic}{écliptique} se déplacerait pourtant de façon inégale en \omterm{def:b1:spherical-astronomy:equatorial}{ascension droite}, la
grandeur qui fixe l'\omterm{def:b1:spherical-astronomy:equatorial}{angle horaire} et donc le cadran solaire~: en retard sur le
Soleil moyen entre un équinoxe et le solstice suivant, en avance entre un
solstice et l'équinoxe suivant. C'est le terme d'obliquité de l'\omterm{def:b1:spherical-astronomy:solar-time}{équation du
temps}.
\end{solution}

\begin{solution}{exo:b1:spherical-astronomy:11}
À son lever, $h = 0$, donc la \cref{prop:b1:spherical-astronomy:transform}
donne $\sin\delta = \cos\varphi\cos\mathrm{A}_0$~; au solstice de juin,
$\delta = \varepsilon$, donc $\cos\mathrm{A}_0 = \sin\varepsilon/\cos\varphi$.
Pour $\mathrm{A}_0 = \ang{45}$~: $\cos\varphi = \sin\ang{23.44}/\cos\ang{45} =
0.5626$, $\varphi = \ang{55.8}$~N. Au-delà de $\ang{90} - \varepsilon =
\ang{66.56}$, $\sin\varepsilon/\cos\varphi > 1$~: il n'y a pas d'\omterm{def:b1:spherical-astronomy:horizontal}{azimut} de
lever, car le Soleil ne se couche pas ce jour-là (Soleil de minuit).
\end{solution}

\begin{solution}{exo:b1:spherical-astronomy:12}
(a) Un fuseau d'angle $\theta$ est la fraction $\theta/2\pi$ de la sphère~:
$\frac{\theta}{2\pi}\times4\pi R^2 = 2\theta R^2$. (b) Les trois paires de
fuseaux ont pour aire totale $2\times(2A + 2B + 2C)R^2$~; elles recouvrent la
sphère une fois, sauf le triangle et son opposé, recouverts chacun trois fois,
soit deux fois de trop~: $4(A + B + C)R^2 = 4\pi R^2 + 4S$, où $S$ est l'aire du
triangle. D'où $S = R^2(A + B + C - \pi)$. (c) Trois angles droits~:
$S = R^2(3\pi/2 - \pi) = \pi R^2/2$, soit un huitième de la sphère (un octant).
Triangle de la figure~: $89.9 + 64.9 + 67.4 - 180 = \ang{42.2}$, c'est-à-dire
\num{0.736} radian~: $S = 0.736\,R^2$, soit \qty{5.9}{\percent} de la sphère.
\end{solution}

\begin{solution}{pb:b1:spherical-astronomy:1}
\textbf{1.} Colatitudes $\ang{90} - (\ang{-24.63}) = \ang{114.63}$ et
$\ang{90} - \ang{51.48} = \ang{38.52}$~; différence de \omterm{def:b1:spherical-astronomy:geographic}{longitude} \ang{70.40}.

\textbf{2.} $\cos\Delta = \cos\ang{114.63}\cos\ang{38.52} +
\sin\ang{114.63}\sin\ang{38.52}\cos\ang{70.40} = -0.3260 + 0.1899 =
-0.1362$.

\textbf{3.} $\Delta = \ang{97.83}$~; distance $6371\times1.7074 =
\qty{10878}{km}$, soit environ \qty{10900}{km}.

\textbf{4.}
\[
  \cos\mathrm{Az} = \frac{\cos\ang{38.52} - \cos\ang{114.63}\cos\ang{97.83}}
  {\sin\ang{114.63}\sin\ang{97.83}} = \frac{0.7823 - 0.0568}{0.9006} = 0.8056 ,
\]
$\mathrm{Az} = \ang{36.3}$, vers le nord-est.

\textbf{5.} $\sin\mathrm{Az} = \sin\ang{70.40}\sin\ang{38.52}/\sin\ang{97.83}
= 0.5922$, ce qui donne \ang{36.3} ou \ang{143.7}~; le cosinus positif de la
question 4 retient \ang{36.3}.

\textbf{6.} Non. Un cap constant coupe tous les méridiens sous le même angle et
décrit une spirale vers un pôle~; le cap du \omterm{def:b1:spherical-astronomy:spherical-triangle}{grand cercle} change continuellement
(il ne vaut \ang{36.3} qu'au départ). L'avion suivrait un trajet plus long et
arriverait ailleurs s'il ne corrigeait pas son cap.

\textbf{7.} La \omterm{def:b1:spherical-astronomy:horizontal}{hauteur} du \omterm{def:b1:spherical-astronomy:celestial-sphere}{pôle céleste} sud est égale à la valeur absolue de la
\omterm{def:b1:spherical-astronomy:geographic}{latitude}~: \ang{24.63}.

\textbf{8.} $h = \ang{90} - |\ang{-24.63} - (\ang{-16.72})| = \ang{82.09}$~;
$\delta > \varphi$, donc Sirius passe \ang{7.91} au nord du \omterm{def:b1:spherical-astronomy:celestial-sphere}{zénith}.

\textbf{9.} Une étoile ne se couche jamais si $\delta < -(\ang{90} - \ang{24.63}) =
\ang{-65.37}$. Canopus, à \ang{-52.70}, n'est pas circumpolaire~: il se couche.

\textbf{10.} $\ang{90} - |\ang{-24.63} + \ang{52.70}| = \ang{61.93}$, au sud
du \omterm{def:b1:spherical-astronomy:celestial-sphere}{zénith}.

\textbf{11.} $\cos H_0 = -\tan(\ang{-24.63})\tan(\ang{-16.72}) =
-(0.4584)(0.3004) = -0.1377$~; $H_0 = \ang{97.92}$. Sirius reste au-dessus de
l'horizon $2\times97.92/15 = \qty{13.06}{h}$ de \omterm{def:b1:spherical-astronomy:sidereal-time}{temps sidéral}, soit environ
\qty{13}{h}.

\textbf{12.} $\cos\mathrm{A}_0 = \sin(\ang{-16.72})/\cos\ang{24.63} =
-0.3165$~: Sirius se lève à l'\omterm{def:b1:spherical-astronomy:horizontal}{azimut} \ang{108.5}, au sud de l'est.

\textbf{13.} Quand $H = 0$, $\mathrm{LST} = \alpha = \qty{6}{h}~\qty{45}{min}~\qty{9}{s}$.

\textbf{14.} $\qty{86400}{s}\times365.2422/366.2422 = \qty{86164.09}{s} =
\qty{23}{h}~\qty{56}{min}~\qty{4}{s}$~; le passage avance de
\qty{235.9}{s}, soit \qty{3}{min}~\qty{56}{s}, chaque nuit.

\textbf{15.} Il faut $12 + n\times24/365.2422 = 6.7525 + 24$, donc
$n = 18.7525/0.065710 = 285.4$ jours. 285 jours après le 20 mars, on est le
30 décembre~: Sirius traverse le méridien à minuit moyen local dans la nuit du
30 au 31 décembre 2026.

\textbf{16.} Le Paranal est à \ang{70.40} à l'ouest de Greenwich~: le temps
moyen local est en retard de $70.40/15 = \qty{4.693}{h}$, soit
\qty{4}{h}~\qty{41.6}{min}.

\textbf{17.} \omterm{def:b1:spherical-astronomy:solar-time}{Temps solaire vrai} = \omterm{def:b1:spherical-astronomy:solar-time}{temps solaire moyen} + \omterm{def:b1:spherical-astronomy:solar-time}{équation du temps} =
$\qty{0}{h} - \qty{3}{min}$~: le cadran solaire indique environ 23~h~57.

\textbf{18.} $7\times\qty{3.93}{min} \approx \qty{27.5}{min}$ plus tôt.

\textbf{19.} $+\ang{23.44}$ au solstice de juin, $-\ang{23.44}$ au solstice de
décembre.

\textbf{20.} Juin~: $\ang{90} - |\ang{-24.63} - \ang{23.44}| = \ang{41.93}$.
Décembre~: $\ang{90} - |\ang{-24.63} + \ang{23.44}| = \ang{88.81}$~; comme
$\delta_\odot = \ang{-23.44} > \varphi$, le Soleil passe \ang{1.19} au nord du
\omterm{def:b1:spherical-astronomy:celestial-sphere}{zénith}.

\textbf{21.} $I = 1361\times(1/0.9837)^2\times\sin\ang{88.81} = 1361\times
1.0334\times0.9998 = \qty{1406}{W.m^{-2}}$.

\textbf{22.} En juin, c'est l'hiver austral. $\cos H_0 = -\tan(\ang{-24.63})
\tan\ang{23.44} = 0.4584\times0.4335 = 0.1988$~: le Soleil se couche à l'\omterm{def:b1:spherical-astronomy:equatorial}{angle
horaire} $H_0 = \ang{78.53}$.

\textbf{23.} $\cos\mathrm{A}_0 = \sin\ang{23.44}/\cos\ang{24.63} = 0.4376$~:
$\mathrm{A}_0 = \ang{64.0}$, au nord de l'est.

\textbf{24.} Durée du jour $2\times78.53/15 = \qty{10.47}{h}$~; la plus longue
nuit de l'année au Paranal dure $24 - 10.47 = \qty{13.53}{h}$, soit
\textbf{\qty[text-series-to-math=true]{13}{h}~\qty[text-series-to-math=true]{32}{min}}.
\end{solution}
